% R32: compute log, memory, and the interruptions (numbers from results_spectra/runlog_v2.txt)
\paragraph{Budget and what was used.} The task came with a budget of 8 hours of compute. The user later lifted it
(``finish within the next 12 hours''). All runs used one 12th-generation i7 with 31~GB of memory, one job at a time.
\begin{itemize}
\item First night (v1 pipeline, 01:53 to 02:51): obc0, all 16 rules for $N=6\ldots13$ and $N=14$ for 9 rules; about
  1.7~h including validation and memory tests.
\item Continuation (v2 pipeline, from 12:32 on 2026-10-14): 485 successful runs and 7.1~h of wall time.
  \begin{center}\small
  \begin{tabular}{p{0.6\textwidth}rr}
  \toprule
  stage & runs & wall time \\
  \midrule
  $H$, obc0 (rest of $N=14$; $N=15\ldots20$ by rule; even-$N$ reruns of the mirror rules with $\varphi$) & 117 & 3.55~h \\
  $H$, rings $N=6\ldots14$ (with $\Omega_T$) & 144 & 2.92~h \\
  $X$, obc0 $N=6\ldots14$ & 144 & 0.44~h \\
  $X$, rings $N=6\ldots14$ & 80 & 0.23~h \\
  \bottomrule
  \end{tabular}
  \end{center}
  Every completed run exited with code 0. Runs cut off by the interruptions below were rerun, except rule~153 at
  $N=15$. The 27 non-zero exit codes in the log are the 12:32 start of the runner, which split paths at the spaces in the repository path; those runs
  failed instantly and were rerun after the fix.
\end{itemize}
The most expensive runs are the single $\approx16\,000$-dimensional sectors at $N=14$: 15 to 18~minutes each, with a peak
resident memory of 14~GB.

\paragraph{Interruptions.}
\begin{enumerate}
\item \emph{First night, 02:51}: the session's memory reaper stopped all jobs. The cause was the dense $U$ of the chaotic
  rules: for these rules $U$ is effectively dense, and storing it as a sparse matrix and running a component search
  over it cost 1.1~GB at $N=12$, or about 18~GB at $N=14$. Fixed in v2 (below); resumed with the user's go-ahead.
\item \emph{About 20:08}: the machine shut down. The run lost was 54 at $N=14$, whose old file had already been removed for
  the rerun with $\varphi$. The user then stopped the runs for a second reboot. The rerun was done after the restart.
\item \emph{About 22:45}: the memory reaper stopped all jobs during rule~153 at $N=15$ (above 17~GB). That run is the
  only one that never finished. It was not restarted (see the next paragraph).
\end{enumerate}

\paragraph{Why the sizes stop where they do.} The pipeline diagonalises each Krylov sector densely (the gate-algebra
blocks are resolved inside it), so the memory is set by the largest \emph{sector}. A $16\,383$-dimensional sector peaks at
about 14~GB, and the memory grows faster than the square of the dimension (201: 2.6~GB at $N=18$, sector 6765; 7.7~GB at
$N=19$, sector $10\,946$). The next sizes would need:
\begin{itemize}
\item 153/195 at $N=15$ and 60/102 at $N=16$: a sector of $2^{15}$, about 50~GB;
\item 150/105 at $N=16$: about twice the $12\,870$-dimensional sector of $N=15$, more than 30~GB;
\item 201 at $N=20$: a sector of $17\,711$ ($F_{22}$), about 20 to 23~GB, which is too close to the limit at which the
  session reaper stops every job.
\end{itemize}
These sizes were therefore not attempted, so every class is compared at the largest $N$ that all its members share. That
is $N=14$ for $k=3$ and for the one-sided rules, $N=15$ for XOR, and $N=19$ for $k=1$ on the open chain; on the ring,
$N=14$ for every rule. Going further would need a block-first diagonalisation, which projects onto the gate-algebra
blocks before the dense step. That is the obvious next optimisation.

\paragraph{The v2 pipeline (\fn{generate\_spectral\_data\_v2.jl}).}
\begin{enumerate}
\item $U$ is never formed: $Uv=L_B(L_Av)$ with the two sparse layers (a few MB at $N=14$).
\item Krylov sectors come from a union-find over the supports of the \emph{individual gates}. For $V=H$ this is the same
  partition as the components of $\mathrm{supp}(U)$ (R-T19 Prop.~4), and it is asserted against the old construction
  for every $N\le10$ run. For $V=X$, $U$ is a permutation and is formed directly.
\item Symmetry checks are random-vector tests through the layers. Dense sector blocks are assembled column-chunk-wise,
  and the products inside each block are computed per batch of at most 2048 columns.
\end{enumerate}
\emph{Validation:} v2 reproduces v1 block by block (sector size, block size, sorted quasienergies, $S_2^\mathrm{Haar}$,
degeneracies) for seven rules at $N=9$. Its ring spectra match the exact Python engine to $<10^{-14}$, for both gates.
Non-degenerate eigenvectors agree to $10^{-15}$ in $S_\mathrm{vN}$; degenerate ones are basis dependent by construction.

\paragraph{Analysis cost.} The first analysis pass looped over blocks with one boolean mask per block, which is
$O(\text{blocks}\times\text{states})$: about $10^{11}$ operations for the $k=1$ rules at $N=20$. It now groups the
states by a single sort (\fn{stats.groups}). The full summary of 173 obc0 units takes 70~s, and its output is identical
to the old one on every unit that was not rerun.
